\documentclass[11pt]{article}
\usepackage[a4paper,margin=22mm]{geometry}
\usepackage{fontspec}
\setmainfont{DejaVu Serif}
\setsansfont{DejaVu Sans}
\usepackage{microtype}
\usepackage{xcolor}
\usepackage{hyperref}
\usepackage{enumitem}
\usepackage{parskip}
\definecolor{Accent}{HTML}{971D32}
\definecolor{Muted}{HTML}{666666}
\hypersetup{colorlinks=true,urlcolor=Accent,linkcolor=Accent}
\setlist{leftmargin=1.6em,itemsep=0.3em,topsep=0.35em}
\setcounter{tocdepth}{2}
\renewcommand{\contentsname}{Contents}
\newcommand{\sectionrule}{\par\vspace{-0.5em}\noindent\textcolor{Accent}{\rule{\linewidth}{0.6pt}}\par}
\begin{document}

\begin{center}
{\sffamily\bfseries\color{Accent} TOEFL WORD OF THE DAY\par}
\vspace{8mm}
{\fontsize{42}{48}\selectfont\bfseries complaisant\par}
\vspace{2mm}
{\Large\sffamily\color{Accent} /kəmˈpleɪzənt/\par}
\vspace{2mm}
{\small\sffamily Local audio phrase: complaisant\par}
{\small\sffamily Audio file: audios/complaisant.mp3\par}
\vspace{2mm}
{\large\sffamily\color{Muted} adjective\par}
\vspace{5mm}
{\sffamily very willing to please or agree\par}
\vspace{6mm}
{\large A complaisant person readily agrees with or accommodates others, sometimes to an excessive degree.\par}
\vspace{6mm}
\begin{minipage}{0.84\textwidth}
\itshape “The complaisant assistant agreed to every request without questioning whether it was reasonable.”
\end{minipage}\par
\vspace{10mm}
\textbf{Date:} August 7th 2026\quad\textbullet\quad \textbf{Level:} B1--mid-B2 TOEFL
\vfill
Created and compiled by \textbf{Amir Kooshky}\par
\vspace{2mm}
\href{https://t.me/KooshkyTOEFL}{Telegram Channel}\quad\textbullet\quad
\href{https://www.instagram.com/kooshkytoefl}{Instagram}\quad\textbullet\quad
\href{https://t.me/KooshkyTOEFL\_pv}{Personal Telegram}
\end{center}
\newpage
\tableofcontents
\newpage

\section{Meanings and usage}
\sectionrule
\subsection{Meaning 1: Very willing to please or agree}
\textbf{Part of speech:} adjective

\textbf{IPA:} /kəmˈpleɪzənt/

\textbf{Pronunciation phrase:} complaisant

\textbf{Local audio file:} audios/complaisant.mp3

\textbf{Register:} formal and relatively uncommon; often disapproving when the willingness to please seems excessive

\textbf{Definition:} very willing to agree with, help, or please other people, sometimes more readily than is wise or sincere

\textbf{In simpler words:} very eager to please or agree

\textbf{Typical pattern:} be complaisant toward + person or group

\subsubsection{Natural examples}
\begin{enumerate}
  \item The complaisant assistant agreed to every request without questioning whether it was reasonable.
  \item He was so complaisant toward senior officials that he rarely expressed his own opinion.
  \item The committee needed an independent adviser, not someone who would remain complaisant in the face of pressure.
  \item Her complaisant manner made difficult customers feel that all their demands would be accepted.
  \item The manager initially appeared helpful, but his complaisant promises were impossible to fulfill.
  \item Critics accused the newspaper of being too complaisant toward the government.
  \item A complaisant negotiator may make unnecessary concessions merely to avoid conflict.
  \item The host was unfailingly complaisant and adjusted the schedule to satisfy every guest.
\end{enumerate}

\subsubsection{Nuance and implication}
\textbf{Central nuance:} The word describes a strong willingness to accommodate other people's wishes or opinions.

\textbf{Usual implication:} It may suggest kindness and cooperation, but it often implies weak judgment, insincerity, or an excessive desire to avoid disagreement.

\textbf{Compared with accommodating:} Accommodating usually describes a helpful willingness to adjust and is generally positive. Complaisant is more formal and often suggests that someone is too willing to please or agree.

\subsubsection{Grammar notes}
\textbf{How to use it:} Use it before a noun, as in a complaisant employee, or after a linking verb, as in The employee was complaisant.

\textbf{What can follow it:} It can be followed by toward and then the person or group whose wishes are being accepted.

\textbf{Common preposition:} Use complaisant toward someone when identifying the person or group being pleased, as in complaisant toward powerful clients.

\textbf{Position in a sentence:} Modifiers such as too, overly, excessively, and remarkably often appear before complaisant.

\subsubsection{Useful patterns}
\textbf{Pattern:} be complaisant toward + person or group

\textbf{Use:} Use this to identify the people whose wishes someone readily accepts.

\textbf{Example:} The regulator was accused of being complaisant toward large corporations.

\textbf{Pattern:} a complaisant + person, manner, attitude, or response

\textbf{Use:} Use this before a noun to describe excessive willingness to please or agree.

\textbf{Example:} His complaisant response failed to address the serious concerns raised by employees.

\textbf{Pattern:} too or overly complaisant

\textbf{Use:} Use this when someone's willingness to please may cause problems.

\textbf{Example:} Leaders who are overly complaisant may avoid making necessary but unpopular decisions.

\textbf{Pattern:} remain complaisant in the face of + pressure or demands

\textbf{Use:} Use this when someone continues to accept what others want despite difficult circumstances.

\textbf{Example:} The board remained complaisant in the face of growing political pressure.

\subsubsection{Common wrong usage}
\paragraph{Correction 1}
\textbf{Wrong:} The company became complaisant after years of success and stopped innovating.

\textbf{Correct:} The company became complacent after years of success and stopped innovating.

\textbf{Why:} Complacent means too satisfied and no longer careful. Complaisant means excessively willing to please or agree with other people.

\paragraph{Correction 2}
\textbf{Wrong:} The employee was complaisant about the safety risk.

\textbf{Correct:} The employee was complacent about the safety risk.

\textbf{Why:} Use complacent about a risk or problem. Complaisant normally describes behavior toward other people.

\paragraph{Correction 3}
\textbf{Wrong:} She complaisanted with every request.

\textbf{Correct:} She was complaisant and agreed to every request.

\textbf{Why:} Complaisant is an adjective, not a verb.

\paragraph{Correction 4}
\textbf{Wrong:} He was complaisant for his demanding clients.

\textbf{Correct:} He was complaisant toward his demanding clients.

\textbf{Why:} Use toward to identify the person or group whose wishes someone readily accepts.

\section{Word family}
\sectionrule
\subsection{complaisance}
\textbf{Part of speech:} uncountable noun

\textbf{IPA:} /kəmˈpleɪzəns/

\textbf{Definition:} a strong willingness to please or agree with other people, especially when it is excessive

\textbf{Relationship to the headword:} This noun names the attitude or behavior described by complaisant.

\textbf{Grammar pattern:} complaisance toward + person or group

\textbf{Usage note:} It is formal and uncommon, and it is normally used without a or an.

\subsubsection{Examples}
\begin{enumerate}
  \item The official's complaisance toward powerful businesses raised concerns about his independence.
  \item Her apparent complaisance concealed considerable disagreement with the proposal.
\end{enumerate}

\subsection{complaisantly}
\textbf{Part of speech:} adverb

\textbf{IPA:} /kəmˈpleɪzəntli/

\textbf{Definition:} in a way that shows great or excessive willingness to please or agree

\textbf{Relationship to the headword:} This adverb describes an action performed in a complaisant manner.

\textbf{Grammar pattern:} complaisantly + agree, accept, smile, or respond

\textbf{Usage note:} It is uncommon and mainly appears in formal or literary writing.

\subsubsection{Examples}
\begin{enumerate}
  \item The assistant complaisantly accepted every change to the schedule.
  \item He smiled complaisantly and assured the clients that all their demands would be met.
\end{enumerate}

\section{Common collocations}
\sectionrule
\subsection{complaisant attitude}
\textbf{Applies to:} Very willing to please or agree

\textbf{Meaning:} an attitude that is excessively willing to accept what others want

\textbf{Example:} The agency's complaisant attitude toward safety violations attracted public criticism.

\subsection{complaisant manner}
\textbf{Applies to:} Very willing to please or agree

\textbf{Meaning:} a polite and highly accommodating way of behaving

\textbf{Example:} His complaisant manner made him popular with demanding customers.

\subsection{complaisant response}
\textbf{Applies to:} Very willing to please or agree

\textbf{Meaning:} a response that readily accepts or supports another person's wishes

\textbf{Example:} The minister gave a complaisant response instead of challenging the inaccurate claim.

\subsection{too complaisant}
\textbf{Applies to:} Very willing to please or agree

\textbf{Meaning:} more willing to please or agree than is appropriate

\textbf{Pattern:} be too complaisant toward + person or group

\textbf{Example:} The board had been too complaisant toward the chief executive.

\subsection{overly complaisant}
\textbf{Applies to:} Very willing to please or agree

\textbf{Meaning:} excessively willing to satisfy other people

\textbf{Example:} An overly complaisant adviser may fail to provide honest criticism.

\subsection{complaisant toward}
\textbf{Applies to:} Very willing to please or agree

\textbf{Meaning:} very ready to accept the wishes or demands of a particular person or group

\textbf{Pattern:} be complaisant toward + person or group

\textbf{Example:} The institution was criticized for being complaisant toward wealthy donors.

\subsection{remain complaisant}
\textbf{Applies to:} Very willing to please or agree

\textbf{Meaning:} continue to accept or accommodate what others want

\textbf{Example:} The representatives remained complaisant despite increasing demands from industry leaders.

\section{Synonyms and antonyms}
\sectionrule
\subsection{Close synonyms}
\subsubsection{accommodating}
\textbf{Part of speech:} adjective

\textbf{Applies to:} Very willing to please or agree

\textbf{Shared meaning:} Both words describe a willingness to adjust to other people's wishes.

\textbf{Difference:} Accommodating normally praises someone for being helpful and flexible. Complaisant is more formal and may criticize someone for being too willing to please.

\textbf{Example:} The hotel staff were extremely accommodating when our plans changed.

\subsubsection{obliging}
\textbf{Part of speech:} adjective

\textbf{Applies to:} Very willing to please or agree

\textbf{Shared meaning:} Both words can describe someone who readily helps or pleases others.

\textbf{Difference:} Obliging usually means pleasantly willing to help and is positive. Complaisant may imply that the willingness is excessive, weak, or insincere.

\textbf{Example:} An obliging neighbor offered to collect our mail.

\subsubsection{agreeable}
\textbf{Part of speech:} adjective

\textbf{Applies to:} Very willing to please or agree

\textbf{Shared meaning:} Both can describe willingness to cooperate with others.

\textbf{Difference:} Agreeable can mean pleasant or willing to accept a suggestion. Complaisant places greater emphasis on a general readiness to satisfy other people.

\textbf{Example:} The committee members were agreeable to changing the meeting time.

\subsubsection{submissive}
\textbf{Part of speech:} adjective

\textbf{Applies to:} Very willing to please or agree

\textbf{Shared meaning:} Both may describe someone who does not strongly resist other people's demands.

\textbf{Difference:} Submissive emphasizes accepting another person's authority or control. Complaisant emphasizes the desire to please or avoid disagreement.

\textbf{Example:} Employees were expected to remain submissive and never question management.

\section{Words learners may confuse}
\sectionrule
\subsection{complacent}
\textbf{Part of speech:} adjective

\textbf{Why learners confuse them:} The words have very similar spellings and pronunciations, but their meanings are different.

\textbf{Key difference:} Complaisant means very or excessively willing to please and agree with others. Complacent means too satisfied with yourself or a situation to notice risks or make improvements.

\textbf{complaisant:} The complaisant adviser agreed with every proposal made by the director.

\textbf{complacent:} The successful company became complacent and stopped preparing for new competitors.

\subsection{compliant}
\textbf{Part of speech:} adjective

\textbf{Why learners confuse them:} Both words can describe cooperation, but compliant focuses on obeying rules or requirements.

\textbf{Key difference:} Complaisant describes a willingness to please people. Compliant means following a rule, request, command, or required standard.

\textbf{complaisant:} The complaisant employee accepted every demand made by the client.

\textbf{compliant:} The equipment is compliant with international safety standards.

\section{Practice exercises}
\sectionrule
\subsection{Word family choice}
Choose the best member of the word family for each blank.

\begin{enumerate}
  \item The official's excessive \_\_\_\_\_ toward powerful companies raised questions about his independence.
  \item The assistant \_\_\_\_\_ accepted every last-minute change.
  \item An adviser should be cooperative without becoming overly \_\_\_\_\_.
\end{enumerate}

\subsection{Correct or incorrect?}
Decide whether each sentence is correct. Correct the inaccurate ones.

\begin{enumerate}
  \item The complaisant employee agreed to every request made by the customer.
  \item The company became complaisant after dominating the market for ten years.
  \item The adviser was too complaisant toward senior officials.
  \item Managers should not be complaisant about cybersecurity risks.
\end{enumerate}

\subsection{Collocation practice}
Complete each sentence with a natural collocation.

\begin{enumerate}
  \item The regulator was criticized for being too complaisant \_\_\_\_\_ large corporations.
  \item His complaisant \_\_\_\_\_ prevented him from questioning unreasonable demands.
  \item An overly \_\_\_\_\_ adviser may tell leaders only what they want to hear.
\end{enumerate}

\subsection{Complete answer key}
\subsubsection{Word family choice}
\begin{enumerate}
  \item \textbf{complaisance.} The official's excessive complaisance toward powerful companies raised questions about his independence. \emph{Use the noun complaisance after excessive.}
  \item \textbf{complaisantly.} The assistant complaisantly accepted every last-minute change. \emph{Use the adverb complaisantly to describe how the assistant accepted the changes.}
  \item \textbf{complaisant.} An adviser should be cooperative without becoming overly complaisant. \emph{Use the adjective complaisant after becoming.}
\end{enumerate}

\subsubsection{Correct or incorrect?}
\begin{enumerate}
  \item \textbf{Correct} \emph{Complaisant correctly describes someone who is highly willing to please or agree.}
  \item \textbf{Incorrect}. The company became complacent after dominating the market for ten years. \emph{Use complacent for excessive satisfaction that reduces effort or caution.}
  \item \textbf{Correct} \emph{Complaisant toward correctly identifies the people whose wishes the adviser accepted too readily.}
  \item \textbf{Incorrect}. Managers should not be complacent about cybersecurity risks. \emph{Use complacent about a risk or problem.}
\end{enumerate}

\subsubsection{Collocation practice}
\begin{enumerate}
  \item \textbf{toward.} The regulator was criticized for being too complaisant toward large corporations. \emph{Use complaisant toward the person or group whose wishes are readily accepted.}
  \item \textbf{attitude.} His complaisant attitude prevented him from questioning unreasonable demands. \emph{A complaisant attitude is an excessively accommodating way of thinking.}
  \item \textbf{complaisant.} An overly complaisant adviser may tell leaders only what they want to hear. \emph{Overly complaisant means excessively willing to please or agree.}
\end{enumerate}

\vfill
\begin{center}
\small\color{Muted} Created and compiled by \textbf{Amir Kooshky}\\
\href{https://t.me/KooshkyTOEFL}{Telegram Channel}\quad\textbullet\quad
\href{https://www.instagram.com/kooshkytoefl}{Instagram}\quad\textbullet\quad
\href{https://t.me/KooshkyTOEFL\_pv}{Personal Telegram}
\end{center}
\end{document}
